% Luc Destombe --- licence CC BY-NC-SA 4.0
% https://creativecommons.org/licenses/by-nc-sa/4.0/deed.fr
% Fichier autonome : compiler avec pdflatex, deux fois (signets, nombre de pages).
\documentclass[a4paper,10pt]{article}
\makeatletter
\newif\iflycee@niveaux
\lycee@niveauxtrue
\newif\iflycee@palatino
\lycee@palatinotrue

\RequirePackage[utf8]{inputenc}
\RequirePackage[T1]{fontenc}
\RequirePackage[french]{babel}
\RequirePackage{etoolbox}

\newcommand{\ShorthandsOff}{\@ifundefined{shorthandoff}{}{\shorthandoff{;:!?}}}
\newcommand{\ShorthandsOn}{\@ifundefined{shorthandon}{}{\shorthandon{;:!?}}}
\ShorthandsOff
\AtBeginDocument{\ShorthandsOn}
\AtBeginEnvironment{tikzpicture}{\ShorthandsOff}

\RequirePackage{geometry}
\RequirePackage{fancyhdr}
\RequirePackage{lastpage}
\RequirePackage{multicol}
\RequirePackage{array}
\RequirePackage{enumitem}
\RequirePackage{xcolor}
\RequirePackage{needspace}

\RequirePackage{amsmath,amssymb}
\RequirePackage{mathpazo}
\DeclareMathSizes{10.95}{10.95}{8}{5.5}
\begingroup\catcode`\_=\active \catcode`\^=\active
\gdef_#1{\sb{\scriptscriptstyle #1}}
\gdef^#1{\sp{\scriptscriptstyle\lycee@moinsepais #1}}
\endgroup
\newcommand\lycee@moins{\mathbin{%
  \setbox\z@\hbox{$\m@th\scriptscriptstyle\lycee@moinsorig$}%
  \dimen@=.55\wd\z@ \advance\dimen@-.5pt
  \hbox to.75\wd\z@{\hss$\m@th\scriptscriptstyle\vcenter{\hbox{%
    \tikz\draw[line width=.5pt,line cap=round](0,0)--(\the\dimen@,0);}}$\hss}}}
\begingroup\lccode`\~=`\- \lowercase{\endgroup\let~}\lycee@moins
\newcommand\lycee@moinsepais{\mathcode`\-="8000 }
\AtBeginDocument{\mathchardef\lycee@moinsorig=\mathcode`\-\relax}
\AtBeginDocument{\catcode`\_=12 \mathcode`\_="8000
  \catcode`\^=12 \mathcode`\^="8000 }
\newcommand\lycee@exposants{%
  \fontdimen13\textfont2=.48\fontdimen6\textfont2
  \fontdimen14\textfont2=.48\fontdimen6\textfont2
  \fontdimen15\textfont2=.48\fontdimen6\textfont2
  \fontdimen13\scriptfont2=.48\fontdimen6\scriptfont2
  \fontdimen14\scriptfont2=.48\fontdimen6\scriptfont2
  \fontdimen15\scriptfont2=.48\fontdimen6\scriptfont2}
\every@math@size\expandafter{\the\every@math@size\lycee@exposants}
\newlength{\retraitcalculs}
\setlength{\retraitcalculs}{2em}

\RequirePackage{tikz}
\usetikzlibrary{arrows.meta,calc,babel,shapes.misc}
\RequirePackage[most]{tcolorbox}
\tcbuselibrary{skins,breakable}

\definecolor{coulDef}{HTML}{1F4E79}
\definecolor{coulProp}{HTML}{7030A0}
\definecolor{coulRem}{HTML}{C55A11}
\definecolor{coulEx}{HTML}{2E7D32}
\definecolor{coulExo}{HTML}{B03A2E}
\definecolor{coulPart}{HTML}{1F4E79}
\definecolor{coulMeth}{HTML}{1B6E4F}
\definecolor{coulCourbe}{HTML}{0B5ED7}
\definecolor{coulCourbeB}{HTML}{C00000}
\definecolor{coulCourbeC}{HTML}{2E7D32}
\definecolor{coulGrille}{HTML}{8C8C8C}
\definecolor{coulRep}{HTML}{1B6E4F}
\definecolor{coulBar}{HTML}{2E7D32}

\newcommand{\R}{\mathbb{R}}
\newcommand{\Rs}{\mathbb{R}^{*}}
\renewcommand{\le}{\leqslant}
\renewcommand{\ge}{\geqslant}
\newcommand{\vect}[1]{\overrightarrow{#1}}
\newcommand{\norme}[1]{\left\lVert #1 \right\rVert}
\newcommand{\intervalle}[2]{\left[#1\,;#2\right]}
\RequirePackage{cancel}
\renewcommand{\CancelColor}{\color{coulExo}}
\newcommand{\fois}[1]{\times{\color{coulExo}#1}}
\newcommand{\barre}[1]{\cancel{#1}}

\tikzset{grille/.style={coulGrille,line width=0.4pt}}
\newcommand{\quadrillage}[4]{\draw[grille,step=1] (#1,#2) grid (#3,#4);}
\tikzset{
  croix/.style={cross out,draw,line width=0.6pt,inner sep=0pt,minimum size=4pt},
  croixdroite/.style={croix,rotate=45,line width=0.8pt},
}
\newcommand{\pt}[3]{\path (#1) node[croix]{} node[#2,font=\small,inner sep=2.5pt]{$#3$};}

\newcommand{\lycee@repere}[4]{%
  \draw[grille,step=1] (#1,#2) grid (#3,#4);
  \draw[-{Stealth[length=2mm]},thick] (#1,0)--({#3+0.4},0) node[below left=-1pt]{$x$};
  \draw[-{Stealth[length=2mm]},thick] (0,#2)--(0,{#4+0.4}) node[below left=-1pt]{$y$};
  \node[below left,font=\scriptsize,inner sep=1.5pt] at (0,0) {$0$};}
\newcommand{\lycee@gradx}[1]{\foreach \k/\l in {#1}{%
  \draw (\k,0.07)--(\k,-0.07) node[below,font=\scriptsize,inner sep=1.5pt]{$\l$};}}
\newcommand{\lycee@grady}[1]{\foreach \k/\l in {#1}{%
  \draw (0.07,\k)--(-0.07,\k) node[left,font=\scriptsize,inner sep=1.5pt]{$\l$};}}
\newcommand{\lycee@intff}[2]{\left[#1\,;#2\right]}
\newcommand{\lycee@intfo}[2]{\left[#1\,;#2\right[}
\newcommand{\lycee@intof}[2]{\left]#1\,;#2\right]}
\newcommand{\lycee@intoo}[2]{\left]#1\,;#2\right[}
\newcommand{\lycee@ens}[1]{\left\{#1\right\}}
\AtBeginDocument{%
  \providecommand{\repere}{\lycee@repere}%
  \providecommand{\gradx}{\lycee@gradx}%
  \providecommand{\grady}{\lycee@grady}%
  \providecommand{\Cf}{\mathcal{C}\sb{\scriptscriptstyle f}}%
  \providecommand{\Cg}{\mathcal{C}\sb{\scriptscriptstyle g}}%
  \providecommand{\intff}{\lycee@intff}%
  \providecommand{\intfo}{\lycee@intfo}%
  \providecommand{\intof}{\lycee@intof}%
  \providecommand{\intoo}{\lycee@intoo}%
  \providecommand{\ens}{\lycee@ens}}

\newlist{questions}{enumerate}{3}
\setlist[questions,1]{label=\arabic*\textdegree),leftmargin=*,itemsep=2pt,topsep=3pt}
\setlist[questions,2]{label=\alph*),leftmargin=*,itemsep=1pt,topsep=2pt}
\setlist[questions,3]{label=\roman*),leftmargin=*,itemsep=1pt,topsep=1pt}
\setlist[itemize]{label=\textbullet,leftmargin=*,itemsep=2pt,topsep=3pt}

\newcommand{\besoinplace}[1]{\par\penalty\@M\begingroup
  \dimen@=#1\relax\dimen@ii=\pagegoal\advance\dimen@ii-\pagetotal
  \ifdim\dimen@>\dimen@ii\ifdim\dimen@ii>\z@\vfil\fi\break\fi\endgroup}

\newcommand{\lycee@pied}{\thepage/\pageref{LastPage}}
\newcommand{\entete}[2]{%
  \pagestyle{fancy}\fancyhf{}%
  \fancyhead[L]{\footnotesize\color{coulPart}\textbf{#1}}%
  \fancyhead[R]{\footnotesize\color{coulPart}\textbf{#2}}%
  \fancyfoot[C]{\footnotesize\lycee@pied}%
  \renewcommand{\headrulewidth}{0.4pt}%
  \renewcommand{\headrule}{\hbox to\headwidth{%
    \color{coulPart!40}\leaders\hrule height \headrulewidth\hfill}}}

\RequirePackage{ccicons}
\newcommand{\lycee@licence}{{\scriptsize\color{gray}%
  \copyright\ Luc Destombe --- \ccbyncsa\ CC BY-NC-SA 4.0}}
\newif\iflycee@licence \lycee@licencetrue
\newcommand{\sanslicence}{\lycee@licencefalse}
\AtBeginDocument{\iflycee@licence\AddToHookNext{shipout/foreground}{%
  \put(0,0){\rlap{\kern\dimexpr1in+\hoffset+\oddsidemargin+\textwidth\relax
    \llap{\raisebox{-\dimexpr1in+\voffset+\topmargin+\headheight+\headsep
      +\textheight+\footskip\relax}{\lycee@licence}}}}}\fi}

\newcommand{\formulesserrees}{%
  \setlength{\abovedisplayskip}{3pt}\setlength{\belowdisplayskip}{3pt}%
  \setlength{\abovedisplayshortskip}{2pt}\setlength{\belowdisplayshortskip}{3pt}}

\setlength{\parindent}{0pt}

\RequirePackage[implicit=false,hidelinks,pdfencoding=auto,psdextra,bookmarksopen,
  bookmarksopenlevel=1,pdfcreator={},pdfproducer={}]{hyperref}
\RequirePackage{bookmark}
\hypersetup{pdfauthor={Luc Destombe},pdfkeywords={CC BY-NC-SA 4.0}}
\newcounter{lycee@signet}
\newcommand{\signet}[3][]{\stepcounter{lycee@signet}%
  \edef\lycee@dest{\if\relax\detokenize{#1}\relax
    signet.\arabic{lycee@signet}\else#1\fi}%
  \hypertarget{\lycee@dest}{}\bookmark[level=#2,dest=\lycee@dest]{#3}}
\pdfstringdefDisableCommands{%
  \def\R{R}\def\vect#1{#1}\def\Cf{Cf}\def\Cg{Cg}\def\fois#1{×#1}%
  \def\barre#1{#1}\def\intervalle#1#2{[#1;#2]}\def\intff#1#2{[#1;#2]}%
  \def\textsuperscript#1{#1}\def\dfrac#1#2{#1/#2}\def\frac#1#2{#1/#2}%
  \def\vec#1{vecteur #1}}
\RequirePackage[official]{eurosym}

\geometry{top=0.8cm,bottom=1.3cm,left=1cm,right=1cm,footskip=0.6cm}

\pagestyle{fancy}
\fancyhf{}
\renewcommand{\headrulewidth}{0pt}
\fancyfoot[C]{\footnotesize Page \thepage{} / \pageref{LastPage}}

\newcommand{\titrefiche}[2]{%
  \begin{center}
    \Large\textbf{#1}\\[2pt]
    \large\textbf{#2}\\[2pt]
  \end{center}
  \vspace{0.10cm}}

\setlength{\columnseprule}{0.4pt}
\setlength{\columnsep}{15pt}
\renewcommand{\columnseprulecolor}{\color{black!45}}
\setlength{\multicolsep}{2pt}

\newcounter{partiefiche}
\renewcommand{\thepartiefiche}{\Roman{partiefiche}}
\newif\ifapresbandeau
\newcommand{\lycee@bandeau}[1]{%
  \par\penalty-600
  \vspace{0.08cm}%
  \noindent\colorbox{coulPart}{%
    \parbox{\dimexpr\linewidth-2\fboxsep\relax}{%
      \color{white}\bfseries\raggedright\strut#1\strut}}%
  \nopagebreak\par\nopagebreak
  \vspace{0.06cm}\nopagebreak
  \apresbandeautrue}
\newcommand{\partiefiche}[1]{%
  \refstepcounter{partiefiche}%
  \lycee@bandeau{\signet{1}{\thepartiefiche{} --- #1}\thepartiefiche{} --- #1}}
\newcommand{\partiefichesansnum}[1]{\lycee@bandeau{\signet{1}{#1}#1}}

\newcounter{exo}
\newlength{\espaceexo}\setlength{\espaceexo}{7pt}
\newlength{\espacetitre}\setlength{\espacetitre}{2pt}
\newcommand{\exo}[1][\relax]{%
  \par
  \ifapresbandeau\apresbandeaufalse\else\penalty-150\addvspace{\espaceexo}\fi
  \refstepcounter{exo}%
  \noindent\signet[exercice-\arabic{exo}]{2}{Exercice \arabic{exo}}%
  \textbf{\color{coulExo}\underline{Exercice \theexo}}%
  \ifx\relax#1\relax\else\textbf{ : #1}\fi
  \par\nopagebreak\vspace{\espacetitre}\nopagebreak
  \ignorespaces}

\setlist[enumerate]{nosep,leftmargin=*,itemsep=2pt,topsep=2pt}
\setlist[itemize]{nosep,leftmargin=*,topsep=2pt}
\newcommand{\choix}[3]{\par\hspace*{1em}%
  \makebox[0.3\linewidth][l]{a) #1}\makebox[0.3\linewidth][l]{b) #2}c) #3}
\newcommand{\deux}[2]{\par\hspace*{0.6em}%
  \makebox[0.48\linewidth][l]{#1}#2}
\newcommand{\trois}[3]{\par\hspace*{0.6em}%
  \makebox[0.32\linewidth][l]{#1}\makebox[0.32\linewidth][l]{#2}#3}

  \renewcommand{\exo}[1][\relax]{%
    \ifapresbandeau\apresbandeaufalse\else\par\penalty-150\fi
    \refstepcounter{exo}%
    \vspace{0.05cm}%
    \noindent\signet[exercice-\arabic{exo}]{2}{Exercice \arabic{exo}}%
    \textbf{\color{coulExo}\underline{Exercice \theexo}}%
    \ifx\relax#1\relax\else\textbf{ : #1}\fi%
    \par\nopagebreak
    \ignorespaces}
  \newcommand{\rappel}[1]{%
    \par\vspace{0.06cm}\nopagebreak
    \noindent\fcolorbox{black!35}{black!5}{%
      \parbox{\dimexpr\linewidth-2\fboxsep-2\fboxrule\relax}{%
        \small\textbf{Rappel.} #1}}%
    \par\vspace{0.06cm}\nopagebreak}
  \setlist[enumerate]{nosep,leftmargin=*}
  \setlist[itemize]{nosep,leftmargin=*}
  \setlength{\multicolsep}{3pt plus 1pt minus 1pt}
  \setlength{\premulticols}{10pt}
  \setlength{\postmulticols}{10pt}
  \newlist{expr}{enumerate}{1}
  \setlist[expr]{label={\Alph*\ $=$},nosep,leftmargin=*,itemsep=0pt}

\renewenvironment{center}{\par\addvspace{3pt}\centering}{\par\addvspace{3pt}}
\formulesserrees
\AtBeginDocument{\formulesserrees}
\clubpenalty=10000
\widowpenalty=10000
\displaywidowpenalty=10000
\brokenpenalty=10000
\setlength{\parskip}{0pt}

\RequirePackage{ragged2e}
\setlength{\RaggedRightRightskip}{0pt plus 1fil}
\AtBeginDocument{\RaggedRight\binoppenalty=10000 \relpenalty=10000 }
\g@addto@macro\@parboxrestore{\RaggedRight}
\makeatother

\tcbuselibrary{listings}
\newtcblisting{python}{%
  listing only,enhanced,
  colback=black!3,colframe=black!35,boxrule=0.5pt,arc=1pt,
  left=4pt,right=4pt,top=1pt,bottom=1pt,before skip=3pt,after skip=3pt,
  listing options={language=Python,basicstyle=\ttfamily\small,
    keywordstyle=\color{coulMeth}\bfseries,columns=fullflexible,
    keepspaces=true,showstringspaces=false,aboveskip=0pt,belowskip=0pt}}
\newcommand{\reperesuite}[4]{%
  \draw[grille,step=1] (-1,#2) grid (#1,#3);
  \draw[-{Stealth[length=2mm]},thick] (-1,0)--({#1+0.4},0)
    node[below left=-1pt,font=\footnotesize]{$n$};
  \draw[-{Stealth[length=2mm]},thick] (0,#2)--(0,{#3+0.4})
    node[below left=-1pt,font=\footnotesize]{$#4$};
  \node[below left,font=\scriptsize,inner sep=1.5pt] at (0,0) {$0$};}

\begin{document}

\titrefiche{1\textsuperscript{re} STI2D --- Suites numériques}{Fiche d'exercices du chapitre}

\begin{multicols}{2}\raggedcolumns

\partiefiche{Définition d'une suite}

\exo[Suites logiques]
Pour chaque liste, dire comment on passe d'un nombre au suivant, puis écrire les trois
nombres suivants.
\begin{multicols}{2}
\begin{enumerate}[label=\alph*)]
  \item $5$ ; $10$ ; $15$ ; $20$ ; $25$ ; \ldots
  \item $2$ ; $9$ ; $16$ ; $23$ ; $30$ ; \ldots
  \item $30$ ; $26$ ; $22$ ; $18$ ; \ldots
  \item $3$ ; $9$ ; $27$ ; $81$ ; \ldots
  \item $1$ ; $8$ ; $27$ ; $64$ ; $125$ ; \ldots
  \item $1$ ; $3$ ; $4$ ; $7$ ; $11$ ; $18$ ; \ldots
\end{enumerate}
\end{multicols}

\exo[Rang et position]
La suite $(u_{n})$ commence par $7$ ; $13$ ; $19$ ; $25$ ; \ldots{}, avec $u_{0}=7$.
\begin{enumerate}[label=\arabic*)]
  \item Donner $u_{1}$ et $u_{3}$.
  \item Quel est le rang du $5^{\text{e}}$ terme ? Calculer ce terme.
  \item Calculer $u_{8}$. À quelle position de la liste se trouve-t-il ?
\end{enumerate}

\exo[Lire un tableau]
Voici les premiers termes d'une suite $(u_{n})$.
\begin{center}
\renewcommand{\arraystretch}{1.25}
\begin{tabular}{|c|c|c|c|c|c|c|}
\hline
$n$ & $0$ & $1$ & $2$ & $3$ & $4$ & $5$\\
\hline
$u_{n}$ & $15$ & $11$ & $7$ & $3$ & $-1$ & $-5$\\
\hline
\end{tabular}
\end{center}
\begin{enumerate}[label=\arabic*)]
  \item Donner $u_{3}$, puis le $3^{\text{e}}$ terme de la suite.
  \item Pour quel rang $n$ a-t-on $u_{n}=-1$ ?
  \item Comment passe-t-on d'un terme au suivant ? En déduire $u_{6}$.
\end{enumerate}

\exo[Les tribunes d'un stade]
Dans une tribune, la première rangée compte $30$ sièges, et chaque rangée en compte $4$ de
plus que la précédente. On note $u_{0}=30$ le nombre de sièges de la première rangée,
$u_{1}$ celui de la deuxième, etc.
\begin{enumerate}[label=\arabic*)]
  \item Calculer $u_{1}$, $u_{2}$, $u_{3}$ et $u_{4}$.
  \item Combien de sièges compte la $10^{\text{e}}$ rangée ?
  \item Quelle rangée compte $70$ sièges ?
\end{enumerate}

\partiefiche{Suites définies explicitement}

\exo[Des entiers]
Soit $u_{n}=5n-8$. Calculer $u_{0}$, $u_{1}$, $u_{2}$, $u_{3}$ et $u_{4}$, puis $u_{20}$ et
$u_{100}$.

\exo[Pour chaque formule]
Pour chaque suite, calculer les quatre premiers termes, puis le terme de rang $10$.
\begin{multicols}{2}
\begin{enumerate}[label=\alph*)]
  \item $a_{n}=3^{n}$
  \item $b_{n}=n^{2}+n$
  \item $c_{n}=-3n+4$
  \item $d_{n}=n^{2}-4n$
  \item $e_{n}=\dfrac{1}{2}n^{2}-1$
  \item $f_{n}=(n-3)^{2}$
\end{enumerate}
\end{multicols}

\exo[Fractions et signes]
\begin{enumerate}[label=\arabic*)]
  \item Soit $u_{n}=\dfrac{n+1}{n+2}$. Calculer $u_{0}$, $u_{1}$, $u_{2}$, $u_{3}$ et
        $u_{8}$, sous forme de fractions.
  \item Soit $v_{n}=(-1)^{n}$. Calculer $v_{1}$, $v_{2}$, $v_{3}$ et $v_{4}$, puis
        $v_{40}$, $v_{65}$ et $v_{100}$.
\end{enumerate}

\exo[$u_{n+1}$ ou $u_{n}+1$ ?]
Pour chaque suite, exprimer $u_{n+1}$, puis $u_{n}+1$, en fonction de $n$.
\begin{multicols}{3}
\begin{enumerate}[label=\alph*)]
  \item $u_{n}=3n+2$
  \item $u_{n}=5-4n$
  \item $u_{n}=n^{2}-2$
\end{enumerate}
\end{multicols}

\exo[La borne de recharge]
Une borne de recharge pour voiture électrique facture $2$~€ la connexion, puis $3$~€ par
heure. On note $p_{n}$ le prix payé pour $n$ heures de charge.
\begin{enumerate}[label=\arabic*)]
  \item Exprimer $p_{n}$ en fonction de $n$.
  \item Combien coûtent $5$ heures de charge ? $8$ heures ?
  \item Une conductrice a payé $23$~€. Combien de temps sa voiture a-t-elle chargé ?
\end{enumerate}

\partiefiche{Suites définies par récurrence}

\exo[Tripler et enlever]
Soit $u_{0}=2$ et $u_{n+1}=3u_{n}-2$. Calculer $u_{1}$, $u_{2}$, $u_{3}$, $u_{4}$ et
$u_{5}$.

\exo[Des relatifs]
Pour chaque suite, calculer les quatre termes qui suivent le premier.
\begin{enumerate}[label=\alph*)]
  \item $u_{0}=4$ et $u_{n+1}=-2u_{n}+3$
  \item $v_{0}=1$ et $v_{n+1}=v_{n}^{2}-3$
  \item $w_{0}=-4$ et $w_{n+1}=6-w_{n}$
\end{enumerate}

\exo[Sans tout calculer ?]
\begin{enumerate}[label=\arabic*)]
  \item Soit $a_{0}=-3$ et $a_{n+1}=-a_{n}$. Calculer $a_{1}$ à $a_{5}$, puis deviner
        $a_{40}$ et $a_{77}$.
  \item Soit $b_{0}=5$ et $b_{n+1}=b_{n}+4$. Calculer $b_{1}$ à $b_{5}$. Combien de
        calculs faudrait-il pour obtenir $b_{50}$ de cette façon ?
\end{enumerate}

\exo[Avec des fractions]
Soit $u_{0}=50$ et $u_{n+1}=\dfrac{10}{u_{n}}$. Calculer $u_{1}$, $u_{2}$, $u_{3}$ et
$u_{4}$. Que remarque-t-on ?

\exo[Quand la relation contient $n$]
Soit $u_{0}=2$ et $u_{n+1}=u_{n}-2n+5$. Calculer $u_{1}$, $u_{2}$, $u_{3}$ et $u_{4}$
(le nombre $n$ change à chaque étape).

\exo[Une culture de levures]
Une cuve contient $300$~milliers de levures. Chaque heure, leur nombre double, puis on en
prélève $80$~milliers. On note $b_{n}$ le nombre de levures (en milliers) après $n$
heures.
\begin{enumerate}[label=\arabic*)]
  \item Donner $b_{0}$, puis exprimer $b_{n+1}$ en fonction de $b_{n}$.
  \item Calculer $b_{1}$, $b_{2}$ et $b_{3}$.
\end{enumerate}

\partiefiche{Programmer une suite définie par récurrence}

\exo[Que renvoie ce programme ?]
\begin{python}
def suite(n):
    u = 3
    for i in range(n):
        u = 3*u - 4
    return u
\end{python}
\begin{enumerate}[label=\arabic*)]
  \item Quelle suite ce programme calcule-t-il ? Donner $u_{0}$ et la relation de
        récurrence.
  \item Calculer à la main ce que renvoient \texttt{suite(1)}, \texttt{suite(2)} et
        \texttt{suite(5)}.
\end{enumerate}

\exo[Même travail]
Même travail avec la ligne \texttt{u = 160} au départ et la ligne \texttt{u = 0.5*u + 2}
dans la boucle, pour \texttt{suite(1)}, \texttt{suite(2)} et \texttt{suite(6)}.

\exo[Compléter le programme]
Soit $u_{0}=2$ et $u_{n+1}=1+2u_{n}$.
\begin{enumerate}[label=\arabic*)]
  \item Calculer $u_{3}$ à la main.
  \item Compléter le programme, puis donner $u_{10}$.
\end{enumerate}
\begin{python}
def suite(n):
    u = .....
    for i in range(n):
        u = .....
    return u
\end{python}
\begin{enumerate}[label=\arabic*),start=3]
  \item Même travail avec $v_{0}=1$ et $v_{n+1}=6-2v_{n}$ : écrire le programme, puis
        donner $v_{1}$, $v_{5}$ et $v_{10}$.
\end{enumerate}

\exo[Les abonnés d'une salle de sport]
Une salle de sport compte $1\,200$ abonnés en 2024. Chaque année, $40\,\%$ des abonnés ne
se réinscrivent pas, et $600$ nouveaux abonnés arrivent. On note $a_{n}$ le nombre
d'abonnés l'année $2024+n$ : $a_{0}=1\,200$.
\begin{enumerate}[label=\arabic*)]
  \item Calculer $a_{1}$ et $a_{2}$.
  \item Exprimer $a_{n+1}$ en fonction de $a_{n}$.
  \item Écrire un programme qui renvoie $a_{n}$.
  \item Combien y aura-t-il d'abonnés en 2039 (arrondi à l'unité) ?
  \item Vers quelle valeur le nombre d'abonnés semble-t-il se stabiliser ?
\end{enumerate}

\partiefiche{Représentation graphique d'une suite}

\exo[Placer les points]
\noindent
\begin{minipage}[t]{0.5\linewidth}\raggedright
Soit $u_{n}=n^{2}-6n+5$.
\begin{enumerate}[label=\arabic*)]
  \item Calculer $u_{0}$, $u_{1}$, \ldots, $u_{6}$.
  \item Représenter ces termes dans le repère ci-contre.
\end{enumerate}
\end{minipage}\hfill
\begin{minipage}[t]{0.47\linewidth}\vspace{0pt}\centering
\begin{tikzpicture}[x=0.42cm,y=0.42cm]
  \reperesuite{7}{-5}{6}{u_{n}}
  \gradx{1,2,3,4,5,6} \grady{-4,-3,-2,-1,1,2,3,4,5}
\end{tikzpicture}
\end{minipage}

\exo[Une suite qui alterne]
Soit $u_{0}=4$ et $u_{n+1}=-u_{n}+2$. Calculer $u_{1}$ à $u_{6}$, puis représenter ces
termes dans un repère.

\exo[Une suite qui se rapproche]
Soit $u_{0}=-6$ et $u_{n+1}=\dfrac{u_{n}+10}{2}$. Calculer $u_{1}$ à $u_{4}$, puis
représenter $u_{0}$ à $u_{4}$ dans un repère. De quel nombre les termes semblent-ils se
rapprocher ?

\exo[Lire un graphique]
\noindent
\begin{minipage}[t]{0.53\linewidth}\raggedright
Les croix représentent les premiers termes d'une suite $(u_{n})$.
\begin{enumerate}[label=\arabic*)]
  \item Lire $u_{0}$, $u_{1}$, $u_{2}$, $u_{3}$ et $u_{4}$.
  \item Comment passe-t-on d'un terme au suivant ?
  \item En déduire $u_{5}$ et $u_{6}$.
\end{enumerate}
\end{minipage}\hfill
\begin{minipage}[t]{0.44\linewidth}\vspace{0pt}\centering
\begin{tikzpicture}[x=0.42cm,y=0.42cm]
  \reperesuite{6}{-1}{9}{u_{n}}
  \gradx{1,2,3,4,5} \grady{1,2,3,4,5,6,7,8}
  \foreach \p in {(0,8),(1,6),(2,4),(3,2),(4,0)} \path[coulCourbe] \p node[croix]{};
\end{tikzpicture}
\end{minipage}

\partiefiche{Sens de variation d'une suite}

\exo[Conjecturer]
Voici les premiers termes de quatre suites. Pour chacune, dire si elle semble
croissante, décroissante, ou ni l'une ni l'autre.
\begin{multicols}{2}
\begin{enumerate}[label=\alph*)]
  \item $1$ ; $3$ ; $6$ ; $10$ ; $15$
  \item $8$ ; $5$ ; $1$ ; $-4$ ; $-10$
  \item $2$ ; $-2$ ; $2$ ; $-2$ ; $2$
  \item $-9$ ; $-7$ ; $-6$ ; $-4$ ; $-3$
\end{enumerate}
\end{multicols}

\exo[Suites définies par récurrence]
Étudier le sens de variation de chaque suite à l'aide de $u_{n+1}-u_{n}$.
\begin{enumerate}[label=\alph*)]
  \item $u_{0}=3$ et $u_{n+1}=u_{n}-6$
  \item $u_{0}=-2$ et $u_{n+1}=u_{n}+n^{2}+1$
  \item $u_{0}=4$ et $u_{n+1}=u_{n}-2n^{2}$
  \item $u_{0}=-5$ et $u_{n+1}=u_{n}+5n+3$
\end{enumerate}

\exo[Suites définies explicitement]
Pour chaque suite, exprimer $u_{n+1}$ en fonction de $n$, puis étudier le sens de
variation à l'aide de $u_{n+1}-u_{n}$.
\begin{multicols}{2}
\begin{enumerate}[label=\alph*)]
  \item $u_{n}=4n-9$
  \item $u_{n}=6-3n$
  \item $u_{n}=2n^{2}-5$
  \item $u_{n}=n^{2}+6n$
\end{enumerate}
\end{multicols}

\exo[Conjecture, puis preuve]
Soit $u_{n}=-n^{2}+4n+5$.
\begin{enumerate}[label=\arabic*)]
  \item Calculer les six premiers termes. Que peut-on conjecturer à partir de $u_{2}$ ?
  \item Montrer que $u_{n+1}-u_{n}=3-2n$.
  \item Quel est le signe de $3-2n$ quand $n\ge 2$ ? Conclure.
\end{enumerate}

\partiefiche{Suites arithmétiques}

\exo[Premiers termes]
$(u_{n})$ est arithmétique. Écrire la relation de récurrence, puis calculer $u_{1}$,
$u_{2}$ et $u_{3}$.
\begin{multicols}{2}
\begin{enumerate}[label=\alph*)]
  \item $u_{0}=4$ et $r=6$
  \item $u_{0}=12$ et $r=-5$
  \item $u_{0}=-3$ et $r=0{,}5$
  \item $u_{0}=\dfrac{1}{4}$ et $r=\dfrac{1}{2}$
\end{enumerate}
\end{multicols}

\exo[Arithmétique ou pas ?]
Pour chaque situation, dire si elle se traduit par une suite arithmétique ; si oui,
donner le premier terme et la raison.
\begin{enumerate}[label=\alph*)]
  \item Un loyer de $650$~€ augmente de $15$~€ chaque année.
  \item Une population de $5\,000$ bactéries augmente de $10\,\%$ par heure.
  \item Une piscine de $60~\text{m}^{3}$ se vide de $4~\text{m}^{3}$ par heure.
  \item Une balle lâchée de $3$~m remonte, à chaque rebond, aux trois quarts de sa
        hauteur précédente.
\end{enumerate}

\exo[Le food truck]
La première semaine, un food truck sert $40$ repas ; ensuite, il en sert chaque semaine
$5$ de plus que la semaine précédente. On note $u_{1}=40$ les repas de la première
semaine, $u_{2}$ ceux de la deuxième, etc.
\begin{enumerate}[label=\arabic*)]
  \item Justifier que $(u_{n})$ est arithmétique et donner sa raison.
  \item Combien de repas sont servis la $2^{\text{e}}$ semaine ? la $6^{\text{e}}$ ?
\end{enumerate}

\exo[Une application]
Une application compte $2\,500$ utilisateurs le premier mois, puis $400$ de plus chaque
mois. On note $v_{1}$ le nombre d'utilisateurs du premier mois, $v_{2}$ celui du
deuxième, etc.
\begin{enumerate}[label=\arabic*)]
  \item Donner $v_{1}$, puis calculer $v_{2}$, $v_{3}$ et $v_{4}$.
  \item Quelle est la nature de la suite $(v_{n})$ ? Donner sa raison.
\end{enumerate}

\partiefiche{Reconnaître une suite arithmétique}

\exo[À partir des termes]
Ces nombres peuvent-ils être les premiers termes d'une suite arithmétique ? Si oui,
donner la raison.
\begin{multicols}{2}
\begin{enumerate}[label=\alph*)]
  \item $3$ ; $8$ ; $13$ ; $18$
  \item $20$ ; $14$ ; $8$ ; $2$
  \item $2$ ; $3$ ; $5$ ; $8$
  \item $1{,}5$ ; $2{,}25$ ; $3$ ; $3{,}75$
  \item $401$ ; $377$ ; $353$ ; $329$
  \item $2$ ; $4$ ; $8$ ; $16$
\end{enumerate}
\end{multicols}

\exo[Démontrer]
Soit $u_{n}=5n-2$.
\begin{enumerate}[label=\arabic*)]
  \item Calculer $u_{0}$, $u_{1}$, $u_{2}$ et $u_{3}$.
  \item Exprimer $u_{n+1}$ en fonction de $n$.
  \item Démontrer que $(u_{n})$ est arithmétique ; préciser son premier terme et sa
        raison.
\end{enumerate}

\exo[Pour chaque suite]
Pour chaque suite, dire si elle est arithmétique. Si oui, donner sa raison et son premier
terme ; sinon, le prouver avec les premiers termes.
\begin{multicols}{2}
\begin{enumerate}[label=\alph*)]
  \item $a_{n}=4n$
  \item $b_{n}=2-3n$
  \item $c_{n}=n^{2}+n$
  \item $d_{n}=\dfrac{n}{3}$
  \item $e_{n}=2-n^{2}$
  \item $f_{n}=(n+2)^{2}-n^{2}$
  \item $g_{0}=1$ et $g_{n+1}=g_{n}^{2}+1$
\end{enumerate}
\end{multicols}

\partiefiche{Terme général et variations d'une suite arithmétique}

\exo[Calculer un terme éloigné]
$(u_{n})$ est arithmétique. Exprimer $u_{n}$ en fonction de $n$, calculer le terme
demandé, puis donner le sens de variation.
\begin{multicols}{2}
\begin{enumerate}[label=\alph*)]
  \item $u_{0}=8$, $r=-3$ : $u_{9}$
  \item $u_{0}=2$, $r=6$ : $u_{100}$
  \item $u_{0}=-5$, $r=\dfrac{5}{2}$ : $u_{6}$
  \item $u_{0}=4$, $r=-\dfrac{2}{3}$ : $u_{9}$
\end{enumerate}
\end{multicols}

\exo[Quand on ne part pas de $u_{0}$]
$(u_{n})$ est arithmétique.
\begin{enumerate}[label=\alph*)]
  \item $u_{1}=3$ et $r=4$. Exprimer $u_{n}$, puis calculer $u_{10}$.
  \item $u_{1}=-6$ et $r=0{,}5$. Calculer $u_{21}$.
  \item $u_{4}=10$ et $r=3$. Calculer $u_{0}$, puis $u_{12}$.
\end{enumerate}

\exo[Une application (suite)]
On reprend l'exercice 31 : $v_{1}=2\,500$ et $r=400$.
\begin{enumerate}[label=\arabic*)]
  \item Exprimer $v_{n}$ en fonction de $n$.
  \item Combien l'application comptera-t-elle d'utilisateurs le $12^{\text{e}}$ mois ?
  \item Quel est le sens de variation de $(v_{n})$ ?
\end{enumerate}

\exo[Économiser pour un scooter]
Une apprentie veut acheter un scooter à $2\,000$~€. Le $1^{\text{er}}$ mars, elle a
$1\,100$~€ ; elle met ensuite $75$~€ de côté chaque mois. On note $u_{1}=1\,100$ son
épargne au $1^{\text{er}}$ mars, $u_{2}$ au $1^{\text{er}}$ avril, etc.
\begin{enumerate}[label=\arabic*)]
  \item Calculer $u_{2}$, $u_{3}$ et $u_{4}$.
  \item Exprimer $u_{n}$ en fonction de $n$, puis calculer $u_{7}$.
  \item Pour quel $n$ a-t-on $u_{n}=2\,000$ ? À quelle date pourra-t-elle acheter le
        scooter ?
\end{enumerate}

\exo[La production d'une usine]
En 2024, une usine fabrique $4\,500$ pièces ; sa production augmente ensuite de $500$
pièces par an. On note $p_{1}=4\,500$ la production de 2024, $p_{2}$ celle de 2025, etc.
\begin{enumerate}[label=\arabic*)]
  \item Exprimer $p_{n}$ en fonction de $n$.
  \item Calculer la production en 2027, puis en 2040.
  \item En quelle année la production aura-t-elle triplé ?
\end{enumerate}

\partiefiche{Suites géométriques}

\exo[Coefficient multiplicateur]
Donner le coefficient multiplicateur $q$ de chaque évolution.
\begin{multicols}{2}
\begin{enumerate}[label=\alph*)]
  \item Augmenter de $4\,\%$
  \item Diminuer de $4\,\%$
  \item Augmenter de $30\,\%$
  \item Diminuer de $30\,\%$
  \item Augmenter de $65\,\%$
  \item Diminuer de $12\,\%$
  \item Augmenter de $200\,\%$
  \item Diminuer de $45\,\%$
\end{enumerate}
\end{multicols}

\exo[Dans l'autre sens]
Pour chaque coefficient multiplicateur, dire s'il traduit une augmentation ou une
diminution, et de quel pourcentage.
\begin{multicols}{4}
\begin{enumerate}[label=\alph*)]
  \item $0{,}96$
  \item $1{,}07$
  \item $0{,}25$
  \item $1{,}6$
  \item $3$
  \item $0{,}98$
  \item $1{,}035$
  \item $0{,}82$
\end{enumerate}
\end{multicols}

\exo[Premiers termes]
$(u_{n})$ est géométrique. Calculer $u_{1}$, $u_{2}$ et $u_{3}$.
\begin{multicols}{2}
\begin{enumerate}[label=\alph*)]
  \item $u_{0}=5$ et $q=3$
  \item $u_{0}=96$ et $q=0{,}5$
  \item $u_{0}=2\,000$ et $q=1{,}05$
  \item $u_{0}=81$ et $q=\dfrac{1}{3}$
\end{enumerate}
\end{multicols}

\exo[Deux situations]
Pour chaque situation, donner le premier terme et la relation entre $u_{n+1}$ et
$u_{n}$, puis calculer $u_{1}$ et $u_{2}$.
\begin{enumerate}[label=\alph*)]
  \item Une technicienne place une prime de $3\,000$~€ à $2\,\%$ par an ; $u_{n}$ est le
        capital après $n$ années.
  \item Une ville de $40\,000$ habitants perd $3\,\%$ de sa population chaque année ;
        $u_{n}$ est la population après $n$ années.
\end{enumerate}

\partiefiche{Reconnaître une suite géométrique}

\exo[À partir des termes]
Ces nombres peuvent-ils être les premiers termes d'une suite géométrique ? Si oui,
donner la raison.
\begin{multicols}{2}
\begin{enumerate}[label=\alph*)]
  \item $5$ ; $10$ ; $20$ ; $40$
  \item $3$ ; $12$ ; $48$ ; $180$
  \item $128$ ; $32$ ; $8$ ; $2$
  \item $4$ ; $8$ ; $12$ ; $16$
  \item $200$ ; $220$ ; $242$ ; $266{,}2$
  \item $50$ ; $10$ ; $2$ ; $0{,}4$
\end{enumerate}
\end{multicols}

\exo[Démontrer]
Soit $u_{n}=2\times 3^{n}$.
\begin{enumerate}[label=\arabic*)]
  \item Calculer $u_{0}$, $u_{1}$, $u_{2}$ et $u_{3}$.
  \item Exprimer $u_{n+1}$ en fonction de $n$.
  \item Démontrer que $(u_{n})$ est géométrique ; préciser son premier terme et sa
        raison.
\end{enumerate}

\exo[Pour chaque suite]
Pour chaque suite, dire si elle est géométrique. Si oui, donner sa raison et son premier
terme ; sinon, le prouver avec les premiers termes.
\begin{multicols}{2}
\begin{enumerate}[label=\alph*)]
  \item $a_{n}=5^{n}$
  \item $b_{n}=n^{2}+3$
  \item $c_{n}=8\times 0{,}5^{n}$
  \item $d_{n}=3n+1$
\end{enumerate}
\end{multicols}
\begin{enumerate}[label=\alph*),start=5]
  \item $e_{0}=1$ et $e_{n+1}=e_{n}^{2}+2$
\end{enumerate}

\exo[La valeur d'une machine]
La valeur d'une machine-outil (en euros) est relevée chaque année :
\begin{center}
\renewcommand{\arraystretch}{1.25}
\begin{tabular}{|c|c|c|c|c|}
\hline
$n$ (années) & $0$ & $1$ & $2$ & $3$\\
\hline
$v_{n}$ & $12\,000$ & $10\,800$ & $9\,720$ & $8\,748$\\
\hline
\end{tabular}
\end{center}
\begin{enumerate}[label=\arabic*)]
  \item Montrer que ces valeurs sont celles d'une suite géométrique, dont on donnera la
        raison.
  \item De quel pourcentage la machine perd-elle de la valeur chaque année ?
  \item Calculer $v_{4}$.
\end{enumerate}

\partiefiche{Terme général et variations d'une suite géométrique}

\exo[Calculer un terme éloigné]
$(u_{n})$ est géométrique. Exprimer $u_{n}$ en fonction de $n$, calculer le terme
demandé, puis donner le sens de variation.
\begin{multicols}{2}
\begin{enumerate}[label=\alph*)]
  \item $u_{0}=3$, $q=2$ : $u_{10}$
  \item $u_{0}=6$, $q=\dfrac{1}{2}$ : $u_{5}$
  \item $u_{0}=2\,000$, $q=0{,}8$ : $u_{5}$
  \item $u_{0}=8$, $q=1$ : $u_{30}$
\end{enumerate}
\end{multicols}

\exo[Quand on ne part pas de $u_{0}$]
$(u_{n})$ est géométrique.
\begin{enumerate}[label=\alph*)]
  \item $u_{1}=4$ et $q=3$. Exprimer $u_{n}$, puis calculer $u_{6}$.
  \item $u_{2}=5$ et $q=2$. Calculer $u_{3}$, $u_{4}$ et $u_{5}$.
\end{enumerate}

\exo[Un placement]
On place $5\,000$~€ à $5\,\%$ par an : chaque fin d'année, les intérêts s'ajoutent au
capital. On note $C_{n}$ le capital après $n$ années.
\begin{enumerate}[label=\arabic*)]
  \item Calculer $C_{1}$ et $C_{2}$.
  \item Quelle est la nature de la suite $(C_{n})$ ? Exprimer $C_{n}$ en fonction de
        $n$.
  \item Calculer le capital après $10$ ans, arrondi au centime.
  \item Au bout de combien d'années le capital aura-t-il doublé ?
\end{enumerate}

\exo[La ville (suite)]
On reprend la ville de l'exercice 43 b) : $u_{0}=40\,000$ et $q=0{,}97$.
\begin{enumerate}[label=\arabic*)]
  \item Exprimer $u_{n}$ en fonction de $n$, puis donner le sens de variation.
  \item Quelle sera la population après $10$ ans ?
  \item Au bout de combien d'années la ville comptera-t-elle moins de $35\,000$
        habitants ?
\end{enumerate}

\exo[La batterie d'un drone]
Au premier vol, la batterie d'un drone tient $30$ minutes. À chaque nouveau vol, elle
s'use : elle ne tient plus que $98\,\%$ de la durée du vol précédent. On note $d_{1}=30$
la durée du premier vol, $d_{2}$ celle du deuxième, etc.
\begin{enumerate}[label=\arabic*)]
  \item Quelle est la nature de la suite $(d_{n})$ ? Exprimer $d_{n}$ en fonction de
        $n$.
  \item Calculer $d_{5}$, arrondi au centième.
  \item À partir de quel vol la batterie tient-elle moins de $24$ minutes ?
\end{enumerate}

\end{multicols}
\end{document}
